\documentclass[../DoAn.tex]{subfiles}
\begin{document}

\begin{center}
    \Large{\textbf{TÓM TẮT NỘI DUNG ĐỒ ÁN}}\\
\end{center}
\vspace{1cm}
Trong bối cảnh chuyển đổi số giáo dục đang diễn ra mạnh mẽ, việc xây dựng các hệ thống quản lý có khả năng tự động hóa và tối ưu hóa vận hành là nhu cầu thiết thực đối với các trung tâm dạy thêm. Đồ án này tập trung vào việc ``Xây dựng hệ thống quản lý trung tâm dạy thêm EduCenter, trọng tâm là xếp lịch thông minh''. Mục tiêu của đề tài là cung cấp nền tảng quản trị vận hành cốt lõi cho trung tâm, đồng thời tạo ra chiều sâu kỹ thuật thông qua bài toán tối ưu xếp lịch học.

Hệ thống đã xây dựng thành công các phân hệ quản trị danh mục gồm học viên, giáo viên, khóa học, chương trình đào tạo, lớp học, phòng học và ca học; đồng thời hình thành chuỗi nghiệp vụ khép kín từ ghi danh, cấu hình lịch tuần, xem trước xếp lịch cho tới sinh các buổi học thực tế. Để giải quyết bài toán xếp lịch, hệ thống chuẩn hóa dữ liệu thời gian bằng thực thể ca học, qua đó bảo đảm các thuật toán cùng vận hành trên một mô hình đầu vào thống nhất.

Điểm nhấn kỹ thuật của đề tài là việc nghiên cứu và áp dụng nhiều thuật toán tối ưu hóa ràng buộc cho cùng một bài toán xếp lịch. Ba hướng tiếp cận đã được cài đặt để so sánh gồm thuật toán tô màu đồ thị kết hợp luật tham lam, thuật toán tìm kiếm Tabu và thuật toán lập trình ràng buộc CP-SAT. Việc đặt các thuật toán trên cùng một bộ dữ liệu đầu vào giúp quá trình đánh giá có cơ sở rõ ràng và hỗ trợ lựa chọn thuật toán chính phù hợp cho hệ thống.

Ngoài phần xếp lịch, đề tài cũng xây dựng trước nền dữ liệu phục vụ học vụ sau buổi học như điểm danh, tổng kết nội dung, kết quả học tập và đơn xin phép. Đây là cơ sở quan trọng để hệ thống có thể phát triển tiếp theo hướng phân tích dữ liệu học tập trong các giai đoạn sau, nhưng không phải là trọng tâm triển khai chính của phiên bản hiện tại.

Các kết quả đạt được bao gồm: hoàn thiện phần lớn các phân hệ quản trị cốt lõi, tổ chức mô hình dữ liệu đủ rộng cho vận hành và mở rộng sau này, xây dựng thành công luồng xem trước xếp lịch, và cài đặt ba thuật toán cho bài toán tối ưu hóa. Hệ thống đã chứng minh được khả năng vận hành thực tế, đồng thời tạo ra chiều sâu học thuật rõ rệt thông qua việc so sánh các chiến lược giải thuật xếp lịch.

\textbf{Từ khóa:} \textit{CP-SAT, Xếp lịch thông minh, Quản lý trung tâm dạy thêm, Tối ưu hóa ràng buộc}

\begin{flushright}
Sinh viên thực hiện\\
\begin{tabular}{@{}c@{}}
\textit{(Ký và ghi rõ họ tên)}
\end{tabular}
\end{flushright}

\end{document}